\section{2023 Probability and Statistics}

Solve every problem.

\subsection{Part I: Probability}

\begin{exercise}\label{probability:2023:1}
Suppose that a sequence $\{X_n\}$ of real-valued random variables converges to $X$ in distribution and there are positive constants $r$ and $C$ such that $\mathbb{E}|X_n|^r \leq C$ for all $n$. Show that
\[
\lim_{n \to \infty} \mathbb{E}|X_n|^s = \mathbb{E}|X|^s
\]
for all $0 < s < r$.
\end{exercise}

\begin{solution}
\textbf{Prerequisites / Core Concepts:}
\begin{itemize}
    \item \textbf{Convergence in Distribution:} A sequence of random variables $X_n$ converges to $X$ in distribution if their cumulative distribution functions converge. By the Continuous Mapping Theorem, if $X_n \xrightarrow{d} X$ and $f$ is a continuous function, then $f(X_n) \xrightarrow{d} f(X)$.
    \item \textbf{Uniform Integrability (UI):} A sequence of random variables $\{Z_n\}$ is uniformly integrable if the "tails" of their expectations shrink to 0 uniformly across all $n$. It is the exact condition needed to upgrade convergence in distribution to convergence in expectation.
    \item \textbf{De la Vallée-Poussin Theorem (Moment bound for UI):} If there exists some $p > 1$ such that $\mathbb{E}[|Z_n|^p]$ is uniformly bounded by a constant $C$, then the sequence $\{Z_n\}$ is uniformly integrable.
\end{itemize}

\textbf{Intuition / Motivation:}

When random variables converge in distribution, their probabilities align, but their expected values don't necessarily converge. Why? Because a tiny probability of an incredibly huge value can artificially inflate the expected value (like a lottery ticket that pays 10 trillion dollars with probability 1 in a trillion).
To guarantee that expectations converge, we need to rule out these "rogue heavy tails." This mathematical safety net is called *Uniform Integrability*.
How do we prove our sequence $|X_n|^s$ doesn't have rogue heavy tails? We look at a higher moment! The problem guarantees that the $r$-th moment is bounded. Since we are interested in the $s$-th moment where $s < r$, the bounded $r$-th moment acts as a massive ceiling that aggressively crushes any probability of the variables taking extremely large values.
By raising our variables to a power slightly larger than 1 (specifically, $p = r/s$), we can use the bounded $r$-th moment to mathematically guarantee uniform integrability. Once uniform integrability is established, the convergence of the expected values follows automatically.

\textbf{Logical Step-by-Step Breakdown:}

\textbf{Step 1: Establish convergence in distribution for the transformed variables.}

We are given that $X_n \xrightarrow{d} X$.
We are interested in the expected value of $|X_n|^s$. Let $f(x) = |x|^s$. Since $s > 0$, the function $f(x)$ is continuous everywhere.
By the Continuous Mapping Theorem, convergence in distribution is preserved under continuous transformations. Therefore:
\[ |X_n|^s \xrightarrow{d} |X|^s \]

\textbf{Step 2: Prove Uniform Integrability using the moment bound.}

To show that $\mathbb{E}[|X_n|^s] \to \mathbb{E}[|X|^s]$, it is sufficient to prove that the sequence of random variables $Z_n = |X_n|^s$ is Uniformly Integrable (UI).
A standard sufficient condition for UI is that the sequence has a uniformly bounded higher moment. We need to find a $p > 1$ such that $\mathbb{E}[Z_n^p] \leq C$ for all $n$.
We are given that $\mathbb{E}[|X_n|^r] \leq C$ for all $n$.
We are also given $0 < s < r$. Let us choose our power to be $p = r/s$. Since $r > s$, we strictly have $p > 1$.
Now we evaluate the $p$-th moment of $Z_n$:
\[ \mathbb{E}[Z_n^p] = \mathbb{E}[(|X_n|^s)^{r/s}] = \mathbb{E}[|X_n|^r] \]
By the given assumption, $\mathbb{E}[|X_n|^r] \leq C$. Thus:
\[ \mathbb{E}[Z_n^p] \leq C \quad \text{for all } n \]
Because the $p$-th moment ($p>1$) is uniformly bounded, the sequence $\{|X_n|^s\}$ is Uniformly Integrable.

\textbf{Step 3: Ensure the limiting variable has a finite expectation.}

We must quickly verify that the target expectation $\mathbb{E}[|X|^s]$ is finite, otherwise the limit would be meaningless.
We use Fatou's Lemma. Since $|X_n|^s \xrightarrow{d} |X|^s$, we can use the Skorokhod representation theorem to assume almost sure convergence on a common probability space without loss of generality.
By Fatou's Lemma:
\[ \mathbb{E}[|X|^s] = \mathbb{E}\left[ \liminf_{n \to \infty} |X_n|^s \right] \leq \liminf_{n \to \infty} \mathbb{E}[|X_n|^s] \]
Since $\mathbb{E}[|X_n|^s]$ is bounded (by Jensen's inequality or Hölder's inequality against the $r$-th moment bound), the right-hand side is finite, meaning $\mathbb{E}[|X|^s] < \infty$.

\textbf{Step 4: Conclude convergence of expectations.}

We have established three critical facts:
1. $|X_n|^s \xrightarrow{d} |X|^s$
2. The sequence $\{|X_n|^s\}$ is Uniformly Integrable.
3. $\mathbb{E}[|X|^s] < \infty$.
By the fundamental theorem of convergence for random variables, if a sequence converges in distribution and is uniformly integrable, then its sequence of expected values converges to the expected value of the limit variable.
Therefore:
\[ \lim_{n \to \infty} \mathbb{E}[|X_n|^s] = \mathbb{E}[|X|^s] \]
\end{solution}

\begin{exercise}\label{probability:2023:2}
Let $p(x,y)$ be the (one-step) transition function of a Markov chain on a discrete state space $S$ and $p_n(x,y)$ be the $n$-step transition function. Show that for any positive integers $L$ and $N$ and any states $x$ and $y$,
\[
\sum_{n=L}^{N+L} p_n(x,y) \leq \sum_{n=0}^N p_n(y,y).
\]
\end{exercise}

\begin{solution}
\textbf{Prerequisites / Core Concepts:}
\begin{itemize}
    \item \textbf{Markov Property:} The future state of a process depends only on the current state, not on the sequence of events that preceded it.
    \item \textbf{Strong Markov Property:} The Markov property holds not just at fixed times, but at *stopping times* (random times whose occurrence is determined solely by the past).
    \item \textbf{First Hitting Time:} The random time step when the chain first visits a specific state. It is a classic example of a stopping time.
\end{itemize}

\textbf{Intuition / Motivation:}

We want to count the expected number of times a particle starting at $x$ visits state $y$ during a specific time window $[L, N+L]$.
Imagine you have a camera that only turns on at time $L$ and turns off at time $N+L$. You want to count how many times the particle visits $y$ while the camera is on.
There are two possibilities:
1. The particle never visits $y$ while the camera is on. The count is 0.
2. The particle visits $y$ at least once. Let's focus on the very *first* time it hits $y$ while the camera is on. Let's call this time $\tau$.
Once the particle hits $y$ at time $\tau$, it's exactly as if it is starting fresh from $y$. The number of *subsequent* visits to $y$ before the camera turns off only depends on how much time is left on the clock.
Because $\tau \geq L$, the maximum amount of time left on the clock is $(N+L) - L = N$.
So, the expected number of visits starting from $x$ is bounded by the expected number of visits if you *started directly at $y$* and had the maximum possible time window (which is $N$ steps). This gives us the strict inequality we are trying to prove.

\textbf{Logical Step-by-Step Breakdown:}

\textbf{Step 1: Write the sum as an expectation of indicator variables.}

Let $X_n$ be the state of the Markov chain at time $n$. The term $p_n(x,y)$ is the probability $\mathbb{P}_x(X_n = y)$, which is identical to the expectation of the indicator function $\mathbb{E}_x[\mathbb{I}(X_n = y)]$.
We can rewrite the left-hand side of the inequality as the expected total number of visits to $y$ in the interval $[L, N+L]$:
\[ \sum_{n=L}^{N+L} p_n(x,y) = \mathbb{E}_x \left[ \sum_{n=L}^{N+L} \mathbb{I}(X_n = y) \right] \]

\textbf{Step 2: Define a stopping time for the first visit.}

Let $\tau_y$ be the first time the chain visits state $y$ during or after time $L$:
\[ \tau_y = \inf\{n \geq L : X_n = y\} \]
$\tau_y$ is a valid stopping time. If the chain never visits $y$ in this window, $\tau_y > N+L$.
We can split the expectation by conditioning on whether the first visit happens within the observation window (i.e., whether $\tau_y \leq N+L$):
\[ \mathbb{E}_x \left[ \sum_{n=L}^{N+L} \mathbb{I}(X_n = y) \right] = \mathbb{E}_x \left[ \mathbb{I}(\tau_y \leq N+L) \sum_{n=\tau_y}^{N+L} \mathbb{I}(X_n = y) \right] \]
(Note that if $\tau_y > N+L$, the sum of indicators inside the window is 0, so we drop that term).

\textbf{Step 3: Apply the Strong Markov Property.}

We condition on the filtration $\mathcal{F}_{\tau_y}$ up to the stopping time $\tau_y$. By the Strong Markov Property, the process $X_{\tau_y + k}$ given $\tau_y$ behaves exactly like a new Markov chain starting from state $X_{\tau_y} = y$.
\[ \mathbb{E}_x \left[ \sum_{n=\tau_y}^{N+L} \mathbb{I}(X_n = y) \mid \mathcal{F}_{\tau_y} \right] = \mathbb{E}_y \left[ \sum_{k=0}^{N+L-\tau_y} \mathbb{I}(X_k = y) \right] \]
where we substituted $k = n - \tau_y$.

\textbf{Step 4: Bound the remaining time window.}

By definition, the stopping time satisfies $\tau_y \geq L$.
Therefore, the number of remaining steps in the window is bounded:
\[ N + L - \tau_y \leq N + L - L = N \]
Because the indicator functions are strictly non-negative, summing them over a shorter time window cannot yield a larger expectation:
\[ \mathbb{E}_y \left[ \sum_{k=0}^{N+L-\tau_y} \mathbb{I}(X_k = y) \right] \leq \mathbb{E}_y \left[ \sum_{k=0}^{N} \mathbb{I}(X_k = y) \right] \]
Notice that the right-hand side is exactly the sum of transition probabilities from $y$ to $y$ over $N$ steps:
\[ \mathbb{E}_y \left[ \sum_{k=0}^{N} \mathbb{I}(X_k = y) \right] = \sum_{n=0}^N p_n(y,y) \]

\textbf{Step 5: Conclude the inequality.}

Substitute the bound back into the overall expectation from Step 2:
\begin{align*}
\sum_{n=L}^{N+L} p_n(x,y) &= \mathbb{E}_x \left[ \mathbb{I}(\tau_y \leq N+L) \mathbb{E}_x \left[ \sum_{n=\tau_y}^{N+L} \mathbb{I}(X_n = y) \mid \mathcal{F}_{\tau_y} \right] \right] \\
&\leq \mathbb{E}_x \left[ \mathbb{I}(\tau_y \leq N+L) \sum_{n=0}^N p_n(y,y) \right]
\end{align*}
Since the sum $\sum_{n=0}^N p_n(y,y)$ is now a deterministic constant (it doesn't depend on the random path), we can pull it out of the expectation:
\[ \sum_{n=L}^{N+L} p_n(x,y) \leq \mathbb{P}_x(\tau_y \leq N+L) \left( \sum_{n=0}^N p_n(y,y) \right) \]
Finally, since probabilities are bounded by 1 (i.e., $\mathbb{P}_x(\tau_y \leq N+L) \leq 1$), we obtain the final strict bound:
\[ \sum_{n=L}^{N+L} p_n(x,y) \leq \sum_{n=0}^N p_n(y,y) \]
\end{solution}

\begin{exercise}\label{probability:2023:3}
Let $\{X_n\}$ be an i.i.d. sequence with the symmetric Bernoulli distribution:
\[
\mathbb{P}\{X = 1\} = \mathbb{P}\{X = -1\} = \frac{1}{2}.
\]
Let $S_n = \sum_{i=1}^n X_i$ be the partial sum. Show that for all $\alpha > \frac{1}{2}$,
\[
\mathbb{P}\left\{\lim_{n \to \infty} \frac{S_n}{n^\alpha} = 0\right\} = 1.
\]
\end{exercise}

\begin{solution}
\textbf{Prerequisites / Core Concepts:}
\begin{itemize}
    \item \textbf{Hoeffding's Inequality:} A powerful concentration inequality that bounds the probability of a sum of bounded independent random variables deviating significantly from its expected value.
    \item \textbf{Borel-Cantelli Lemma:} If the sum of the probabilities of a sequence of events is finite, then the probability that infinitely many of them occur is 0. This is the standard tool for proving almost sure convergence.
    \item \textbf{Almost Sure Convergence:} $Z_n \to 0$ almost surely if, for any $\epsilon > 0$, the probability of $|Z_n| > \epsilon$ occurring infinitely often is zero.
\end{itemize}

\textbf{Intuition / Motivation:}

We are summing random coin flips ($+1$ or $-1$). The expected sum is 0. By the Central Limit Theorem, the sum $S_n$ typically grows at a rate proportional to $\sqrt{n}$ (or $n^{0.5}$).
The problem asks us to divide $S_n$ by $n^\alpha$ where $\alpha > 0.5$. Because we are dividing by something that grows strictly faster than the typical size of $S_n$ (which is $n^{0.5}$), the fraction $S_n / n^\alpha$ should naturally shrink to 0.
But we need to prove this happens *almost surely*, meaning there's a $100\%$ chance that the sequence permanently settles down to 0, despite the occasional crazy string of luck.
To prove this, we ask: what is the probability that $S_n / n^\alpha$ randomly jumps above some tiny threshold $\epsilon$? Because we are dividing by a large denominator, $S_n$ would have to be huge (greater than $\epsilon n^\alpha$). Because $\alpha > 0.5$, this requirement grows much faster than the standard deviation. Hoeffding's inequality tells us that the probability of such an extreme deviation decays exponentially fast. The exponential decay is so rapid that the sum of these probabilities is finite, allowing the Borel-Cantelli lemma to lock the sequence down to 0.

\textbf{Logical Step-by-Step Breakdown:}

\textbf{Step 1: Set up the threshold for almost sure convergence.}

To prove $\lim_{n \to \infty} \frac{S_n}{n^\alpha} = 0$ almost surely, it is sufficient to show that for any arbitrarily small constant $\epsilon > 0$, the event $E_n = \left\{ \left| \frac{S_n}{n^\alpha} \right| > \epsilon \right\}$ occurs only finitely many times.
Equivalently, we want to bound the probability:
\[ \mathbb{P}(|S_n| > \epsilon n^\alpha) \]

\textbf{Step 2: Bound the probability using Hoeffding's Inequality.}

The random variables $X_i$ are independent and strictly bounded within the interval $[-1, 1]$. Their mean is $\mathbb{E}[X_i] = 0$.
Hoeffding's inequality states that for the sum $S_n = \sum X_i$ of variables bounded by $[a, b]$, the probability of deviating from the mean by more than $t$ is:
\[ \mathbb{P}(|S_n - \mathbb{E}[S_n]| > t) \leq 2 \exp\left( - \frac{2t^2}{\sum_{i=1}^n (b-a)^2} \right) \]
Here, $a = -1$, $b = 1$, so $(b-a)^2 = 2^2 = 4$. Our threshold is $t = \epsilon n^\alpha$.
Substitute these into the inequality:
\[ \mathbb{P}(|S_n| > \epsilon n^\alpha) \leq 2 \exp\left( - \frac{2 (\epsilon n^\alpha)^2}{n \cdot 4} \right) = 2 \exp\left( - \frac{2 \epsilon^2 n^{2\alpha}}{4n} \right) \]
Simplifying the exponent:
\[ \mathbb{P}(|S_n| > \epsilon n^\alpha) \leq 2 \exp\left( - \frac{\epsilon^2}{2} n^{2\alpha - 1} \right) \]

\textbf{Step 3: Analyze the exponent and sum the probabilities.}

Look closely at the power of $n$ in the exponential: $2\alpha - 1$.
We are given that $\alpha > 1/2$. Therefore, $2\alpha > 1$, which means $2\alpha - 1 = \delta > 0$ for some strictly positive $\delta$.
Our probability bound is:
\[ \mathbb{P}(E_n) \leq 2 \exp\left( - \frac{\epsilon^2}{2} n^\delta \right) \]
Because $\delta > 0$, the term $n^\delta$ grows to infinity, making the exponential term decay to 0 extremely rapidly (faster than any polynomial).
We must check if the infinite sum of these probabilities converges:
\[ \sum_{n=1}^\infty \mathbb{P}(E_n) \leq \sum_{n=1}^\infty 2 \exp\left( - \left(\frac{\epsilon^2}{2}\right) n^\delta \right) \]
Since exponential decay with any positive fractional power beats polynomial growth, this series converges to a finite number.

\textbf{Step 4: Apply the Borel-Cantelli Lemma.}

The first Borel-Cantelli Lemma states that if the sum of the probabilities of a sequence of events is finite ($\sum \mathbb{P}(E_n) < \infty$), then the probability that infinitely many of them occur is strictly 0.
\[ \mathbb{P}(E_n \text{ infinitely often}) = 0 \]
Because this holds for any $\epsilon > 0$, the sequence $\frac{S_n}{n^\alpha}$ must eventually stay within $(-\epsilon, \epsilon)$ for all large $n$. Thus, it converges to 0 almost surely:
\[ \mathbb{P}\left\{\lim_{n \to \infty} \frac{S_n}{n^\alpha} = 0\right\} = 1 \]
\end{solution}

\begin{exercise}\label{probability:2023:4}
Let $X^n = \{X_{ij}\}$ be an $n \times n$ random matrix whose entries are i.i.d. with symmetric Bernoulli distribution:
\[
\mathbb{P}\{X = 0\} = \mathbb{P}\{X = 1\} = \frac{1}{2}.
\]
Let $p_n = \mathbb{P}\{\det X^n \text{ is odd}\}$. Show that $\lim_{n \to \infty} p_n > 0$.
\end{exercise}

\begin{solution}
\textbf{Prerequisites / Core Concepts:}
\begin{itemize}
    \item \textbf{Modulo Arithmetic for Determinants:} The determinant of a matrix composed of integers is odd if and only if its determinant evaluated modulo 2 is non-zero (i.e., $\det(X) \equiv 1 \pmod 2$).
    \item \textbf{Finite Fields ($\mathbb{F}_2$):} A number system with only two elements $\{0, 1\}$. Addition is XOR ($1+1=0$). A matrix over $\mathbb{F}_2$ is invertible if and only if its determinant is non-zero (which means it must be 1).
    \item \textbf{Linear Independence over Finite Fields:} Just like in standard linear algebra, a matrix is invertible if and only if its rows form a basis (are linearly independent). The span of $k$ independent vectors in $\mathbb{F}_2^n$ contains exactly $2^k$ vectors.
\end{itemize}

\textbf{Intuition / Motivation:}

We want to know the probability that a giant random matrix of 0s and 1s has an odd determinant.
Thinking about the actual integer value of the determinant is a nightmare because it involves adding and subtracting millions of products.
However, we only care about parity (even vs. odd). We can do all the math modulo 2 right from the start! In modulo 2, the determinant is either 0 (even) or 1 (odd). A determinant of 1 means the matrix is invertible in the world of modulo-2 arithmetic.
How do we build an invertible matrix row by row?
- The first row just can't be all zeros.
- The second row can't be all zeros, AND it can't be a copy of the first row.
- The third row can't be any linear combination (using 0s and 1s) of the first two rows.
Because we are in a discrete world, we can just count the number of "bad" combinations at each step and divide by the total number of possible rows ($2^n$). We get a product of probabilities. As the matrix gets infinitely large, does this probability shrink to zero? Surprisingly, no! The chance of randomly picking a linearly dependent row drops exponentially as we add more rows, meaning the infinite product stabilizes to a strictly positive constant.

\textbf{Logical Step-by-Step Breakdown:}

\textbf{Step 1: Convert the problem to Linear Algebra over $\mathbb{F}_2$.}

Let $X^n$ be the $n \times n$ matrix. We map its entries to the finite field $\mathbb{F}_2$.
Since $\det(X^n)$ is a polynomial in the matrix entries, reducing the entries modulo 2 gives the determinant modulo 2.
The determinant is odd if and only if $\det(X^n) \pmod 2 \neq 0$, which in $\mathbb{F}_2$ means $\det(X^n) = 1$.
This is equivalent to saying the matrix $X^n$ is invertible over $\mathbb{F}_2$.
Thus, $p_n$ is simply the probability that a random $n \times n$ matrix over $\mathbb{F}_2$ is invertible.

\textbf{Step 2: Construct the invertible matrix row by row.}

A matrix is invertible over a field if and only if its $n$ rows are linearly independent. We can calculate the probability by selecting rows sequentially.
There are $2^n$ total possible vectors for each row.
- **Row 1:** It must be linearly independent from the empty set. This just means it cannot be the zero vector.
There is exactly $1$ zero vector, so there are $2^n - 1$ valid choices.
The probability of a valid 1st row is $\frac{2^n - 1}{2^n}$.
- **Row 2:** It must be linearly independent from Row 1. This means it cannot be in the span of Row 1.
The span of a single non-zero vector in $\mathbb{F}_2$ contains exactly $2^1 = 2$ vectors (the zero vector and the vector itself).
So there are $2^n - 2$ valid choices.
The probability of a valid 2nd row is $\frac{2^n - 2}{2^n}$.
- **Row $k$:** It must be linearly independent from the first $k-1$ rows. Assuming the first $k-1$ rows are independent, their span is a subspace of dimension $k-1$.
A subspace of dimension $k-1$ over $\mathbb{F}_2$ contains exactly $2^{k-1}$ vectors.
Therefore, any of these $2^{k-1}$ vectors is a "bad" choice. There are $2^n - 2^{k-1}$ valid choices.
The probability of a valid $k$-th row, given the previous rows are independent, is $\frac{2^n - 2^{k-1}}{2^n}$.

\textbf{Step 3: Calculate the overall probability $p_n$.}

Since the choices are made independently at random, we multiply the probabilities for all $n$ rows:
\[ p_n = \prod_{k=1}^n \frac{2^n - 2^{k-1}}{2^n} \]
We can simplify each fraction by dividing the numerator and denominator by $2^n$:
\[ \frac{2^n - 2^{k-1}}{2^n} = 1 - 2^{k-1-n} = 1 - 2^{-(n-k+1)} \]
So, $p_n = \prod_{k=1}^n (1 - 2^{-(n-k+1)})$.
Let's reverse the order of the product by substituting $j = n - k + 1$. As $k$ goes from $1$ to $n$, $j$ goes from $n$ down to $1$.
\[ p_n = \prod_{j=1}^n (1 - 2^{-j}) \]

\textbf{Step 4: Analyze the limit as $n \to \infty$.}

We want to show that $\lim_{n \to \infty} p_n > 0$.
The limit is an infinite product: $P = \prod_{j=1}^\infty (1 - 2^{-j})$.
For an infinite product $\prod (1 - a_j)$ with $0 \leq a_j < 1$ to converge to a strictly positive number, it is a standard theorem in analysis that it is necessary and sufficient for the sum $\sum a_j$ to converge.
In our case, $a_j = 2^{-j}$. The sum is a geometric series:
\[ \sum_{j=1}^\infty 2^{-j} = \frac{1/2}{1 - 1/2} = 1 < \infty \]
Since the sum converges, the infinite product converges to a strictly positive limit.
(Numerically, this limit is approximately $0.288788$).
Therefore, $\lim_{n \to \infty} p_n > 0$.
\end{solution}

\subsection{Part II: Statistics}

\begin{exercise}\label{probability:2023:5}
You have been asked to help design a randomized trial of a new drug, drug A, to be used in place of drug B, for a particular medical condition. The budget is fixed to have 1000 patients treated with A and 1000 treated with drug B. Describe a class of methods that achieves this goal where each patient has a positive probability of receiving drug A and a positive probability of receiving drug B. Provide enough detail that you are describing an explicit algorithm.
\end{exercise}

\begin{solution}
\textbf{Prerequisites / Core Concepts:}
\begin{itemize}
    \item \textbf{Randomized Controlled Trial (RCT):} A study design where treatments are randomly assigned to subjects, ensuring that the treatment groups are statistically comparable, minimizing confounding bias.
    \item \textbf{Random Allocation Rule (Complete Randomization):} A specific randomization algorithm that guarantees exact predetermined sample sizes for each treatment group while maintaining an equal probability of assignment across the entire population.
\end{itemize}

\textbf{Intuition / Motivation:}

We need to allocate exactly 1000 patients to drug A and 1000 to drug B.
If we just flip a coin for every patient, we might end up with 1020 getting A and 980 getting B. That violates the strict budget constraint.
If we assign the first 1000 to A and the next 1000 to B, we satisfy the budget, but we completely destroy randomization (patients arriving late might be sicker, older, etc., creating massive bias).
The solution is to "shuffle a deck of cards." Imagine a deck with 1000 red cards (Drug A) and 1000 black cards (Drug B). We shuffle the deck perfectly. When a patient arrives, we draw the top card.
This guarantees exactly 1000 of each will be handed out (the deck size is fixed). It also guarantees that before the drawing begins, every single patient's chance of drawing a red card is exactly 50\% (because the deck is perfectly shuffled).

\textbf{Logical Step-by-Step Breakdown:}

\textbf{Step 1: Identify the constraints.}
1. Exact allocation: Total A must be 1000, Total B must be 1000.
2. Randomness: Every patient must have a probability $\mathbb{P}(Z_i = A) > 0$ and $\mathbb{P}(Z_i = B) > 0$.

\textbf{Step 2: Propose the Random Allocation Rule (Urn Model / Shuffling).}
The most standard algorithm in clinical trial design to achieve this is the **Random Allocation Rule**.
\textit{Algorithm Description:}
1. Create a fixed array or list $V$ of length $N = 2000$.
2. Populate the first 1000 elements with the label 'A' and the remaining 1000 elements with the label 'B'.
3. Generate a uniformly random permutation of this array. This can be done programmatically using the Fisher-Yates (Knuth) shuffle algorithm.
4. As the $i$-th patient enters the trial (for $i = 1, \dots, 2000$), assign them the treatment corresponding to the $i$-th element of the shuffled array.

\textbf{Step 3: Verify the budget constraint.}
By construction, the array $V$ contains exactly 1000 'A's and 1000 'B's. The permutation step only changes their order, not their counts. Therefore, when all 2000 patients have been processed, exactly 1000 will have received A and 1000 will have received B.

\textbf{Step 4: Verify the positive probability constraint.}
Consider the $i$-th patient to arrive. What is their marginal probability of receiving drug A before the trial begins?
Because the permutation is chosen uniformly at random from all $2000!$ possible permutations, the $i$-th position in the array is equally likely to be occupied by any of the original 2000 elements.
Since 1000 of those elements are 'A', the marginal probability is:
\[ \mathbb{P}(Z_i = A) = \frac{1000}{2000} = 0.5 > 0 \]
Similarly, $\mathbb{P}(Z_i = B) = 0.5 > 0$.
The problem only requires *marginal* probabilities to be positive. (Note: Conditional probabilities also remain positive during the trial as long as there are still 'A's and 'B's remaining in the conceptual urn, but the marginal unconditional probability is what formally defines the fairness of the design from the outset).
\end{solution}

\begin{exercise}\label{probability:2023:6}
You are given the results of a randomized experiment of two drugs, A and B. The experiment was not conducted in the usual way, but rather by allocating patients by a machine-learning algorithm under which each patient has a positive probability of receiving A and of receiving B; moreover the algorithm is completely specified and is built to create better than random balance on the covariates.

(a) Can unbiased estimates of the causal effect of drug A versus B be found, and if so, show why?

(b) Can exact small sample, non-parametric inferences for the causal effect in part (a) be derived, based solely on the randomization distribution of some statistic?
\end{exercise}

\begin{solution}
\textbf{Prerequisites / Core Concepts:}
\begin{itemize}
    \item \textbf{Causal Inference and Potential Outcomes:} Every patient has two potential outcomes: $Y_i(1)$ if given drug A, and $Y_i(0)$ if given drug B. The causal effect is $Y_i(1) - Y_i(0)$.
    \item \textbf{Inverse Probability Weighting (IPW):} An estimation technique that weights observations by the inverse of their probability of receiving the treatment they actually received. This creates a "pseudo-population" free of treatment assignment bias.
    \item \textbf{Randomization Test / Fisher's Exact Test:} A non-parametric method for testing a sharp null hypothesis by calculating a test statistic for every possible permutation of treatment assignments that *could have* occurred under the experimental design.
\end{itemize}

\textbf{Intuition / Motivation:}

Instead of flipping a coin, an AI assigned treatments. But the AI didn't just pick deterministically; it rolled weighted dice for each patient based on their traits (e.g., an older person might have a 70\% chance of getting A, a younger person 40\%).
(a) Can we get an unbiased estimate? Yes. If you know *exactly* how the AI weighted the dice, you can correct for it! If older people were pushed toward drug A (70\% chance), there will be too many older people in group A. To fix this, we take the older people who *did* get A and down-weight them (multiply by 1/0.7). We take the rare older people who got B (30\% chance) and up-weight them (multiply by 1/0.3). By doing this "Inverse Probability Weighting" across the board, we perfectly balance the scales and eliminate the AI's bias, allowing an unbiased estimate.
(b) Can we test if the drug works without assuming normal distributions or large samples? Yes. If the drug does absolutely nothing (the "sharp null"), then your outcome was already predetermined, regardless of what the AI assigned you. We can ask the AI: "If we ran this experiment 10,000 times, what are all the different ways you might have assigned the drugs?" For each hypothetical assignment, we calculate the difference in means. This creates a custom, exact distribution of our statistic based purely on the AI's known programming. We then see where our *actual* experiment's result falls on this custom distribution to get an exact p-value.

\textbf{Logical Step-by-Step Breakdown:}

\textbf{Step 1: Unbiased Estimation using IPW (Part a).}

Let $Z_i \in \{0, 1\}$ be the treatment indicator (1 for A, 0 for B). Let the potential outcomes be $Y_i(1)$ and $Y_i(0)$. The observed outcome is $Y_i = Z_i Y_i(1) + (1-Z_i)Y_i(0)$.
The ML algorithm provides a known probability of assignment for each patient, conditional on their covariates $X_i$:
\[ \pi_i = \mathbb{P}(Z_i = 1 \mid X_i) \]
The problem states that $\pi_i \in (0, 1)$ for all patients (positive probability for both A and B).
We can construct the Inverse Probability Weighting (IPW) estimator for the Average Treatment Effect (ATE):
\[ \hat{\tau}_{IPW} = \frac{1}{N} \sum_{i=1}^N \left( \frac{Z_i Y_i}{\pi_i} - \frac{(1 - Z_i) Y_i}{1 - \pi_i} \right) \]

\textbf{Step 2: Prove the estimator is unbiased.}

We take the expectation of the estimator over the randomization distribution (the random assignment $Z_i$, given the fixed covariates and potential outcomes).
Focus on the first term:
\[ \mathbb{E}\left[ \frac{Z_i Y_i}{\pi_i} \right] = \mathbb{E}\left[ \frac{Z_i Y_i(1)}{\pi_i} \right] \]
Since $Y_i(1)$ and $\pi_i$ are fixed/determined by covariates, we only take the expectation of $Z_i$:
\[ \mathbb{E}[Z_i] = \mathbb{P}(Z_i = 1) = \pi_i \]
Therefore:
\[ \mathbb{E}\left[ \frac{Z_i Y_i(1)}{\pi_i} \right] = \frac{\pi_i Y_i(1)}{\pi_i} = Y_i(1) \]
By exactly symmetric logic for the second term, $\mathbb{E}\left[ \frac{(1 - Z_i) Y_i}{1 - \pi_i} \right] = Y_i(0)$.
Thus, $\mathbb{E}[\hat{\tau}_{IPW}] = \frac{1}{N} \sum (Y_i(1) - Y_i(0)) = ATE$. The estimate is unbiased.

\textbf{Step 3: Exact Small Sample Inference (Part b).}

To perform exact inference, we formulate Fisher's Sharp Null Hypothesis:
$H_0$: The treatment has absolutely no effect on any unit. Mathematically, $Y_i(1) = Y_i(0)$ for all $i$.
Under $H_0$, the potential outcomes are identical to the observed outcomes $Y_i$ regardless of the treatment assignment. The only source of randomness in the entire experiment is the treatment assignment vector $\mathbf{Z} = (Z_1, \dots, Z_N)$.

\textbf{Step 4: Construct the Randomization Distribution.}

Because the ML algorithm is "completely specified," we know the exact joint probability distribution $\mathbb{P}(\mathbf{Z} = \mathbf{z})$ for any possible treatment vector $\mathbf{z} \in \{0, 1\}^N$.
Choose a test statistic that measures treatment effect, for example, the simple difference in means $T(\mathbf{Z}, \mathbf{Y_{obs}})$.
Under $H_0$, $\mathbf{Y_{obs}}$ is fixed. We can re-calculate $T(\mathbf{z}, \mathbf{Y_{obs}})$ for every single possible assignment vector $\mathbf{z}$ that the ML algorithm could have generated.
By plotting these values of $T$ against the known probability $\mathbb{P}(\mathbf{Z} = \mathbf{z})$, we generate the exact reference distribution of the test statistic under the null hypothesis.
We then locate our actually observed test statistic $T_{obs}$ on this distribution to calculate an exact, finite-sample p-value. This process makes zero assumptions about population distributions, making it entirely non-parametric. Thus, exact inference is possible.
\end{solution}
